\section*{Sammanfattning}
\addcontentsline{toc}{section}{Sammanfattning}

\begin{itemize}
  \item En seriell länk levererar bytes och ingenting annat. Var ett meddelande börjar, hur långt
    det är, vad det betyder, om det kom fram helt och vem det kom från är frågor strömmen inte
    besvarar.
  \item En \textbf{frame} är en avgränsad, självbeskrivande enhet som besvarar dem. Kursens frame
    har åtta fält: SOF, LEN, TYPE, DST, SRC, SEQ, DATA och CHK.
  \item \textbf{SOF} är två bytes därför att en byte förr eller senare dyker upp i nyttolasten.
    SOF ensamt räcker ändå inte --- längd och checksumma måste också stämma.
  \item \textbf{LEN} behövs därför att mottagaren läser en byte i taget och inte kan se framåt.
  \item \textbf{Checksumman} ger integritet, inte leveransgaranti. Den säger att framen kom fram
    hel, inte att den kom fram.
  \item Flerbytesfält skrivs \textbf{big-endian}, byte för byte med explicita skiftningar ---
    aldrig genom att kopiera minnet.
\end{itemize}
